\documentclass[10pt,a4paper]{amsart}
\usepackage[margin=19mm]{geometry}
\usepackage{amsmath,amssymb,mathtools}
\usepackage[scr=rsfs,cal=euler]{mathalfa}
\usepackage{xfrac}
\usepackage[unicode,hidelinks,bookmarks=false]{hyperref}
\newcommand{\ie}{i.e.\ }
\newcommand{\cf}{cf.\ }
\newcommand{\resp}{respectively\ }
\newcommand{\Subsection}{\subsection}
\providecommand{\nobreakdash}{\nobreak\hbox{-}}
\setlength{\emergencystretch}{2em}
\multlinegap0pt
\usepackage{mathtools}
\usepackage{amssymb}
\usepackage{xfrac}
\usepackage[shortlabels]{enumitem}
\newtheorem{theorem}{Theorem}
\newcommand{\E}{{\mathbb{E}}}
\newcommand{\R}{\mathbb{R}}
\newcommand{\ibias}{\operatorname{IBias}^2}
\DeclareMathOperator{\ibr}{IBr}
\DeclareMathOperator{\ivar}{IVar}
\DeclareMathOperator{\mise}{MISE}
\title{Bias Reduction by Multiscale Quasi-Interpolation for Scalar- and Manifold-Valued Functions}
\author{Asaf Abas, Nir Sharon}
\date{Source: arXiv:2507.06707v4 (CC BY 4.0)}
\begin{document}
\maketitle

\noindent\emph{Source and licence.} Bias Reduction by Multiscale Quasi-Interpolation for Scalar- and Manifold-Valued Functions. Version arXiv:2507.06707v4 (CC BY 4.0), distributed by arXiv under the Creative Commons Attribution 4.0 International licence, http://creativecommons.org/licenses/by/4.0/, as recorded in the arXiv licence metadata for this exact version and shown on the abstract page https://arxiv.org/abs/2507.06707v4. Full licence notice: \texttt{licenses/arxiv-2507.06707-CC-BY-4.0.txt}.\par\smallskip

\noindent\textit{Editorial scope.} Counted unit: the article's theorem thm:spectral-damping (source0.tex lines 486--514). The article's complete original proof of it is delivered with it (source0.tex lines 516--560, one \texttt{proof} environment of 45 lines). No mathematics is written, repaired or re-derived: the statement and the proof are the article's own lines, in the article's own notation, and no retained line is substituted or re-flowed.  Retained uncounted as the source-local context the proof needs: lines 218--218; lines 267--292.  Nothing was omitted from any retained span, so the package carries no omission record.  No retained line contains a citation, so the wrapper carries no bibliography and no bibliography entry.  No retained line contains a figure environment, an image-inclusion command or drawing code, so no image is delivered and the package declares no asset dependency.  Editorial formatting only, in addition: an AMS \texttt{amsart} preamble that restates the article's own declarations and macros used by the retained lines, namely the environments \texttt{theorem} on separate counters, together with the standard packages those lines need.  The retained environments print in this document as Theorem 1 in order of appearance, whereas the article prints the same items with the numbering of its own full document.  The complete native source of the article as distributed in the arXiv e-print for this exact version is bundled unchanged as \texttt{sources/181283/source0.tex}.
\subsection{Multiscale with linear quasi-interpolation operators}\label{sec:linear}

\begin{theorem}[Bias equals the deterministic multiscale residual on the data sites]\label{thm:bias-equals-residual}
All quantities below are taken on the data-site space~$\mathcal H=\R^{X}$ introduced at the beginning of Section~\ref{sec:linear}. Let~$Y=f+\varepsilon$ be the noisy observation of~$f\in\mathcal H$, where~$\E[\varepsilon]=0$. Let~$\widetilde f_n=M_nY$ be the scalar-valued multiscale estimator after~$n$ levels, with the hierarchy fixed as above. Then,
\begin{equation}\label{eq:expected-multiscale-estimator}
    \E[\widetilde f_n]=M_nf,
\end{equation}
and the conditional bias is
\begin{equation}\label{eq:bias-residual-identity}
    f-\E[\widetilde f_n]=E_nf=\mathcal R_nf.
\end{equation}
Consequently, at any data site~$x \in X$, the conditional squared bias is
\begin{equation}\label{eq:pointwise-bias-residual}
   \left(f(x)-\E[\widetilde f_n(x)]\right)^2=\left((\mathcal R_nf)(x)\right)^2.
\end{equation}

If, in addition,~$\operatorname{Cov}(\varepsilon)=\sigma^2 I$, then the conditional integrated quantities satisfy
\begin{equation}\label{eq:integrated-bias--variance-residual}
\begin{split}  
    \ibias_n &= \|\mathcal R_nf\|^2,
    \quad
    \ivar_n = \sigma^2\operatorname{tr}(M_nM_n^*),
     \\
    \mise_n &= \|\mathcal R_nf\|^2+\sigma^2\operatorname{tr}(M_nM_n^*),
    \end{split}
\end{equation}
where~$M_n^*$ is the adjoint of~$M_n$.
\end{theorem}

\begin{theorem}[Spectral damping of bias]
\label{thm:spectral-damping}
Assume the setting of Theorem~\ref{thm:bias-equals-residual}, with \(\operatorname{Cov}(\varepsilon)=\sigma^2 I\). Suppose that \(Q_1,\ldots,Q_n\) share an orthonormal eigenbasis \(\{u_k\}_{k=1}^N\) and satisfy
\[
    Q_i u_k=\lambda_{i,k}u_k,
    \qquad
    0\leq \lambda_{i,k}\leq 1 .
\]
Let \(f=\sum_{k=1}^N\theta_k u_k\), and define
\[
    a_{j,k}:=\prod_{i=1}^j(1-\lambda_{i,k}),
    \qquad j=1,\ldots,n .
\]
Then
\[
    \ibias_j
    =
    \sum_{k=1}^N a_{j,k}^2\theta_k^2,
    \qquad
    \ivar_j
    =
    \sigma^2\sum_{k=1}^N(1-a_{j,k})^2 .
\]
Define
\[
    \ibr_j:=\frac{\ibias_j}{\ibias_j+\ivar_j}
\]
whenever the denominator is nonzero. Then~\(\ibias_j\) is nonincreasing in~\(j\), \(\ivar_j\) is nondecreasing in~\(j\), and~\(\ibr_j\) is nonincreasing in~\(j\).
\end{theorem}

\begin{proof}
Since \(Q_i u_k=\lambda_{i,k}u_k\), the residual operator \(R_i=I-Q_i\) satisfies
\[
    R_i u_k=(1-\lambda_{i,k})u_k .
\]
Therefore, for the first \(j\) multiscale levels,
\[
    \mathcal R_j u_k
    =
    R_jR_{j-1}\cdots R_1u_k
    =
    a_{j,k}u_k .
\]
Hence
\[
    \mathcal R_j f
    =
    \sum_{k=1}^N a_{j,k}\theta_k u_k .
\]
By Theorem~\ref{thm:bias-equals-residual},
\[
    \ibias_j=\|\mathcal R_j f\|^2
    =
    \sum_{k=1}^N a_{j,k}^2\theta_k^2 .
\]
Since \(M_j=I-\mathcal R_j\), we also have
\[
    M_j u_k=(1-a_{j,k})u_k .
\]
Using the variance formula from Theorem~\ref{thm:bias-equals-residual},
\[
    \ivar_j
    =
    \sigma^2\operatorname{tr}(M_jM_j^*)
    =
    \sigma^2
    \sum_{k=1}^N (1-a_{j,k})^2 .
\]

Finally, because \(0\leq \lambda_{i,k}\leq 1\),
\[
    0\leq a_{j+1,k}\leq a_{j,k}\leq 1 .
\]
Thus \(a_{j,k}^2\) is nonincreasing in \(j\), while \((1-a_{j,k})^2\) is nondecreasing in \(j\). Therefore \(\ibias_j\) is nonincreasing and \(\ivar_j\) is nondecreasing. Finally, for nonnegative \(s\) and \(t\) with \(s+t>0\), the quantity \(s/(s+t)\) is nondecreasing in \(s\) and nonincreasing in \(t\); applying this with \(s=\ibias_j\) and \(t=\ivar_j\) proves the monotonicity of~\(\ibr_j\).
\end{proof}

\end{document}
